\documentclass[11pt]{article}
\usepackage[margin=25mm]{geometry}
\usepackage{fontspec}
\setmainfont{Latin Modern Roman}
\usepackage{amsmath,amssymb,mathtools}
\usepackage{unicode-math}
\setmathfont{Latin Modern Math}
\usepackage{xurl,hyperref}
\hypersetup{colorlinks=true,urlcolor=blue}
\setcounter{secnumdepth}{0}
\setlength{\parindent}{0pt}
\setlength{\parskip}{5pt}
\setlength{\emergencystretch}{3em}
\newcommand{\sref}[2]{\hyperlink{source:#1:#2}{[S#2]}}
\newcommand{\cataloguescope}[1]{\paragraph{Recorded target.}#1}
\newcommand{\R}{\mathbb R}
\newcommand{\C}{\mathbb C}
\newcommand{\N}{\mathbb N}
\newcommand{\Z}{\mathbb Z}
\newcommand{\Q}{\mathbb Q}
\newcommand{\F}{\mathbb F}
\newcommand{\eps}{\varepsilon}
\begin{document}
% UnsolvedMath ID 30002833; source catalogue code OWR-13673-001
\setcounter{section}{30002832}
\section{Conformal Field Theory Realizations of Subfactors}\label{30002833}
{\small UnsolvedMath ID 30002833 · Mathematical Physics}
\subsection{Definitions and mathematical statement}

\paragraph{Shared notation.}$\mathbb N=\{1,2,\ldots\}$; the standard number fields have their usual meanings. All additional conventions are stated locally.


\paragraph{Definitions and conventions.}
A subfactor is an inclusion $N\subset M$ of von Neumann algebras with trivial centers. Its index and standard invariant encode the size and higher relative-commutant structure of the inclusion. In an operator-algebraic conformal field theory, interval or local observable algebras and sectors can produce subfactor inclusions. To say a subfactor arises in conformal field theory requires identifying the inclusion or invariant that the construction realizes. The primary report concerns finite Jones index, with the intended conformal realization recovering the standard invariant or at least its quantum double. This finite-index context is explicit in Vaughan Jones\textquotesingle s contribution, even though the catalogue sentence abbreviates it.

\paragraph{Target retained from the canonical source catalogue.}
\paragraph{Catalogue wording retained for traceability.}
Does every subfactor arise in conformal field theory?

\paragraph{Source scope and qualification.}
The primary contribution (PDF page 5, printed page 853) supplies the finite-index and realization qualifications omitted from the abbreviated catalogue sentence. ID 30004278 revisits the same broad construction programme, and ID 30002840 is the Asaeda--Haagerup instance; equivalence of every realization convention has not been established.
\subsection{Short English statement}
Can every subfactor be realized by a conformal field theory?
\subsection{Research attempt}
\paragraph{Model, date and investigated scope.}
OpenAI GPT-6 Astra (configured model \texttt{gpt-6-astra}), ultra reasoning; generated 11 October 2026. Provides an explicit equal-dimension different-fusion-data test, ruling out scalar index/dimension matching as a sufficient realization check.

\paragraph{Sources and literature check.}
The complete supplied definition, including source qualifications and the dated literature assessment, was reviewed. The calculation below is an elementary restricted attempt motivated by [S2]; it does not claim novelty or a fresh literature-status determination.

\paragraph{Context and investigated scope.} The complete supplied definition, including source qualifications and the dated literature assessment, was reviewed. The calculation below is an elementary restricted attempt motivated by [S2]; it does not claim novelty or a fresh literature-status determination.

\paragraph{Why index alone cannot specify a realization.} Even at the algebraic level, recovering a standard invariant requires more than reproducing a scalar size. For a concrete illustration, take pointed fusion rules on a finite group $G$: simple symbols $X_g$ multiply by $X_gX_h=X_{gh}$. Both groups $\mathbb Z/4\mathbb Z$ and $(\mathbb Z/2\mathbb Z)^2$ have four symbols of dimension one and the same total squared dimension four, but their multiplication tables are not isomorphic. The first has an element of order four, whereas every nonidentity element in the second squares to the identity.

\paragraph{What this tests.} If a proposed conformal realization checks only a matching total dimension, it cannot distinguish even these simple tensor-data possibilities. An equivalence respecting simple objects and products preserves element orders in this pointed case, so the difference is intrinsic. This is a calculation about candidate invariant data, not an assertion that the two examples give a particular index-four subfactor or that their conformal realizations are unavailable.

\paragraph{Source-level gap.} The source concerns finite-index subfactors and asks for the actual inclusion/standard invariant or at least its quantum double. A positive-energy local conformal net must additionally supply interval algebras, covariance and locality, and its sectors must be matched to the required invariant. None is constructed from arbitrary subfactor data here. The first missing step is such a reconstruction with a verified tensor/braiding comparison, not merely an index match. The Asaeda--Haagerup instance 30002840 and the repeated broad program 30004278 remain within that distinction.

\paragraph{Verification and adversarial audit.}
Checked pointed multiplication tables and element-order invariant; made no unproved subfactor realization assertion from the toy data. Model self-audit only; no independent human proof checking. Prompts A and B were applied to this canonical record and its qualified source scope. Checked pointed multiplication tables and element-order invariant; made no unproved subfactor realization assertion from the toy data. Retained unresolved status and zero original-target progress; A local positive-energy conformal reconstruction matching the full standard invariant or quantum double is missing.

Coordinator checked the distinct pointed fusion multiplication tables despite equal total squared dimension. This additional check was model-assisted, not independent human verification.

\paragraph{Tools, process and contributions.}
Python 3 was used for file preparation and structural checks; no numerical experiment or theorem prover was used. Contribution prompt A was followed by prompt B. GPT-6 Astra with ultra reasoning drafted the research and performed its stated self-audit; any additional coordinator contribution is identified in the verification and audit. A GPT-6-based Codex coordinator assembled the manuscript and checked repository integration; its exact variant and reasoning configuration were not exposed. No human mathematical contribution, independent human proof check or formal proof certification is supplied. Complete derivations and the consulted source references are included here; private reasoning transcripts are not published.

\paragraph{Final status and first remaining gap.}
\textbf{UNRESOLVED}. A local positive-energy conformal reconstruction matching the full standard invariant or quantum double is missing.

\section{Completion Estimate}
\noindent\textbf{COMPLETION: 0\%}
The calculation is an auxiliary or restricted attempt and establishes no newly verified portion of the full source target. A local positive-energy conformal reconstruction matching the full standard invariant or quantum double is missing.

\subsection{Sources}
\begin{itemize}
\item[\hypertarget{source:30002833:1}{S1}] ULAM.AI and the UnsolvedMath Contributors, UnsolvedMath, record 30002833 (OWR-13673-001). Catalogue curation under CC BY 4.0; definitions and English exposition added locally.
\url{https://huggingface.co/datasets/ulamai/UnsolvedMath}.
\item[\hypertarget{source:30002833:2}{S2}] Subfactors and Conformal Field Theory (2015), Oberwolfach Reports, DOI 10.4171/owr/2015/16, pp. 849--926. Primary report accessed 2026-10-11; relevant PDF pages 5.
\url{https://doi.org/10.4171/owr/2015/16}.
\item[\hypertarget{source:30002833:3}{S3}] Bischoff--Davydov--Kawahigashi--Longo, index below four (catalogue-listed context; not independently consulted in this preparation).
\url{https://arxiv.org/abs/1506.02606}.
\end{itemize}
\label{end:30002833}
\end{document}
